%! suppress = Makeatletter
%! suppress = TooLargeSection
%! suppress = MissingLabel
\documentclass{article}

% Fields
\usepackage{geometry}
\geometry{top=25mm}
\geometry{bottom=35mm}
\geometry{left=20mm}
\geometry{right=20mm}
% ------------------------------------------------

% Graphics
\usepackage{color}
\usepackage{tabularx}
\usepackage{tikz}
\usepackage{blkarray}
\usepackage{graphicx}
% ------------------------------------------------

% Math
\usepackage{amsmath, amsfonts}
\usepackage{amssymb}
\usepackage{proof}
\usepackage{mathrsfs}
% Crossed-out symbols
% https://tex.stackexchange.com/questions/75525/how-to-write-crossed-out-math-in-latex
\usepackage[makeroom]{cancel}
\usepackage{mathtools}
% ------------------------------------------------

% Additional font sizes
% https://www.overleaf.com/learn/latex/Questions/How_do_I_adjust_the_font_size%3F
\usepackage{moresize}
% Additional colors
% https://www.overleaf.com/learn/latex/Using_colours_in_LaTeX
\usepackage{xcolor}
% \texttimes
\usepackage{textcomp}
% ------------------------------------------------

% Language
\usepackage[utf8] {inputenc}
\usepackage[T2A] {fontenc}
\usepackage[english, russian] {babel}
\usepackage{indentfirst, verbatim}
\usetikzlibrary{cd, babel}
% ------------------------------------------------

% Fonts
\usepackage{stmaryrd}
\usepackage{cmbright}
\usepackage{wasysym}
% ------------------------------------------------

% Code
% https://tex.stackexchange.com/questions/99475/how-to-invoke-latex-with-the-shell-escape-flag-in-texstudio-former-texmakerx
% Colored, requires --shell-escape compiling option
\usepackage{minted}
\setminted{xleftmargin=\parindent, autogobble, escapeinside=\#\#}
\usepackage{listings}
% ------------------------------------------------

% Custom envs
% https://tex.stackexchange.com/questions/371286/draw-a-horizontal-line-in-latex
\newenvironment{proof}{\subparagraph{\hspace{-1em}Решение:\newline}\-}{\par\noindent\rule{\textwidth}{0.4pt}}
% ------------------------------------------------

% Custom commands
\newcommand{\comb}[1]{\mathbf{#1}}
\newcommand{\step}{\rightsquigarrow}
\newcommand{\term}[1]{\mathbf{#1}}
\newcommand{\ap}{~}
\newcommand{\termdef}{\coloneqq}
\newcommand{\subst}[3]{\left[#2 \mapsto #3 \right] #1}
\newcommand{\eqbeta}{=_\beta}
\newcommand{\betastep}{\rightarrow_\beta}
\newcommand{\betasteps}{\twoheadrightarrow_\beta}
\newcommand{\eqeta}{=_\eta}
\newcommand{\tlist}[1]{\term{[}#1\term{]}} % list-term
\renewcommand{\emph}[1]{{\color{blue} \textit{#1}}}
% ------------------------------------------------

% Head
\usepackage{fancyhdr}
\usepackage{hyperref}
\pagestyle{fancy}
\fancyhead[R]{Ярошевский Илья}
\fancyhead[L]{ИТМО MSE, ФП 2024, Дз 1}
% ------------------------------------------------

% Numbering
% https://tex.stackexchange.com/questions/80113/hide-section-numbers-but-keep-numbering
\makeatletter
\renewcommand\thesubsection{Блок \@arabic\c@subsection.\hspace{-0.8em}}
\renewcommand\thesubsubsection{Задание \@arabic\c@subsection.\@arabic\c@subsubsection\hspace{-0.8em}}
% https://tex.stackexchange.com/questions/327689/numbering-subsubsections-with-letters
\renewcommand\theparagraph{\alph{paragraph})\hspace{-0.8em}}
% https://tex.stackexchange.com/questions/129208/numbering-paragraphs-in-latex
\setcounter{secnumdepth}{4}
\makeatother
% ------------------------------------------------

\begin{document}
    \section*{Дз 1. Лямбда-исчисление}

    Глобальный минимум: 8б.

    \subsection{Синтаксис и семантика лямбда-исчисления}

    Локальный минимум: 4б.

    \subsubsection{(1б)}

    Уберите все лишние скобки из выражений.

    Подсказка: можно воспользоваться представлением терма как дерева.

    \paragraph{}

    $(\lambda x\ldotp (x~(x)))$

    \begin{proof}
        \(\lambda x\ldotp x~x\)
    \end{proof}

    \paragraph{}

    $(\lambda x~y \ldotp f~((g~x)~y))$

    \begin{proof}
        \(\lambda x~y~f\ldotp f~(g~x~y)\)
    \end{proof}

    \paragraph{}

    $x ~ (\lambda y\ldotp (\lambda z \ldotp x~y~z))$

    \begin{proof}
        x~\lambda y~z\ldotp x~y~z
    \end{proof}

    \subsubsection{(1б)}

    Найдите все свободные переменные, а также обозначьте, к какому связывателю относится каждая связанная (это удобнее всего сделать с помощью индексов).

    \paragraph{}

    $\lambda x \ldotp \lambda x \ldotp x~y$

    \begin{proof}
         Свободные: \(y\)

        \(\lambda x^0\ldotp \lambda x^1\ldotp x^1~y\)
    \end{proof}

    \paragraph{}

    $\lambda z \ldotp (\lambda x \ldotp x)~(x~z)$

    \begin{proof}
         Свободные: \(x\)

        \(\lambda z^0 \ldotp (\lambda x^1 \ldotp x^1)~(x~z^0)\)
    \end{proof}

    \paragraph{}

    $(\lambda x\ldotp y~y~(\lambda y\ldotp x~((\lambda z\ldotp z~y~x)~(\lambda x\ldotp z~x))))$

    \begin{proof}
         Свободные: \(y\)

        \((\lambda x^0\ldotp y~y~(\lambda y^1\ldotp x^0~((\lambda z^2\ldotp z^2~y^1~x^0)~(\lambda x^3\ldotp z^2~x^3))))\)
    \end{proof}

    \subsubsection{}

    \paragraph{(1б)}

    Покажите, что $\comb{K}^* \eqbeta \comb{K}~\comb{I}$.

    \begin{proof}
        \begin{align}
            \comb{K}~\comb{I} & \equiv (\lambda x~y\ldotp x) \lambda x\ldotp x \betastep \\
                             & \betastep \lambda y\ldotp \lambda x\ldotp x =_\alpha \\
                             & =_\alpha \lambda x~y\ldotp y \equiv \comb{K}^*
        \end{align} 
    \end{proof}

    \paragraph{(1б)}

    Покажите, что $\comb{B} \eqbeta \comb{S}~(\comb{K}~\comb{S})~\comb{K}$, расписывая и нумеруя промежуточные шаги.

    Подсказка: не спешите сразу раскрывать все определения.

    \begin{proof}
        \begin{align}
            \comb{S}~(\comb{K}~\comb{S})~\comb{K} & \equiv (\lambda f~g~x\ldotp f~x~(g~x))~(\comb{K}~\comb{S})~\comb{K} \betastep \\
                             & \betastep (\lambda g~x\ldotp (\comb{K}~\comb{S})~x~(g~x))~\comb{K} \betastep \\
                             & \betastep \lambda x\ldotp (\comb{K}~\comb{S})~x~(\comb{K}~x) \equiv \\
                             & \equiv \lambda x\ldotp ((\lambda x~y\ldotp x)~\comb{S})~x~((\lambda x~y\ldotp x)~x) \betasteps \\
                             & \betasteps \lambda x\ldotp (\lambda y\ldotp \comb{S})~x~(\lambda y\ldotp x) \betastep \\
                             & \betastep \lambda x\ldotp \comb{S}~(\lambda y\ldotp x) \equiv \\
                             & \equiv \lambda x^0\ldotp (\lambda f~g~x^1\ldotp f~x^1~(g~x^1))~(\lambda y\ldotp x^0) \betastep \\
                             & \betastep \lambda x^0\ldotp \lambda g~x^1\ldotp (\lambda y\ldotp x^0)~x^1~(g~x^1) \betastep \\
                             & \betastep \lambda x^0\ldotp \lambda g~x^1\ldotp x^0~(g~x^1) =_\alpha \\
                             & =_\alpha \lambda f~g~x\ldotp f~(g~x) \equiv \comb{B}
        \end{align} 
    \end{proof}

    \subsubsection{}

    Осуществите подстановку и приведите к нормальной форме~--- форме, в которой не осталось редексов, указывая при этом основные шаги.

    \paragraph{(1б)}

    $x \termdef \comb{S}$ в $x~(\lambda z~x\ldotp z~x)~(\lambda z~y~x\ldotp x)~x$.

    \begin{proof}
        \begin{align}
            & \comb{S}~(\lambda z~x\ldotp z~x)~(\lambda z~y~x\ldotp x)~\comb{S} \equiv \\
            \equiv~ & (\lambda f~g~x\ldotp f~x~(g~x))~(\lambda z~x\ldotp z~x)~(\lambda z~y~x\ldotp x)~\comb{S} \betasteps \\
            \betasteps~ & (\lambda z~x\ldotp z~x)~\comb{S}~((\lambda z~y~x\ldotp x)~\comb{S}) \betasteps \\
            \betasteps~ & \comb{S}~((\lambda z~y~x\ldotp x)~\comb{S}) \equiv \\
            \equiv~ & (\lambda f~g~x\ldotp f~x~(g~x))~((\lambda z~y~x\ldotp x)~\comb{S}) \betasteps \\
            \betasteps~ & \lambda g~x\ldotp ((\lambda z~y~x\ldotp x)~\comb{S})~x~(g~x) \betasteps \\
            \betasteps~ & \lambda g~x\ldotp g~x
        \end{align} 
    \end{proof}

    \paragraph{(1б)}

    $x \termdef \comb{K}$ и $y \termdef \comb{B}$ в $y~(\comb{K}~(\lambda z\ldotp x~x~(x~z~x))~y)~(\lambda x\ldotp x~x~x)~\comb{I}$.

    \begin{proof}
        \begin{align}
            & \comb{B}~(\comb{K}~(\lambda z\ldotp \comb{K}~\comb{K}~(\comb{K}~z~\comb{K}))~\comb{B})~(\lambda x\ldotp x~x~x)~\comb{I} \equiv \\
            \equiv~ & (\lambda f~g~x\ldotp f~(g~x))~(\comb{K}~(\lambda z\ldotp \comb{K}~\comb{K}~(\comb{K}~z~\comb{K}))~\comb{B})~(\lambda x\ldotp x~x~x)~\comb{I} \betasteps \\
            \betasteps~ & (\comb{K}~(\lambda z\ldotp \comb{K}~\comb{K}~(\comb{K}~z~\comb{K}))~\comb{B})~((\lambda x\ldotp x~x~x)~\comb{I}) \equiv \\
            \equiv~ & ((\lambda x~y\ldotp x)~(\lambda z\ldotp \comb{K}~\comb{K}~(\comb{K}~z~\comb{K}))~\comb{B})~((\lambda x\ldotp x~x~x)~\comb{I}) \betasteps \\
            \betasteps~ & (\lambda z\ldotp \comb{K}~\comb{K}~(\comb{K}~z~\comb{K}))~((\lambda x\ldotp x~x~x)~\comb{I}) \betasteps~ \\
            \betasteps~ & \comb{K}~\comb{K}~(\comb{K}~((\lambda x\ldotp x~x~x)~\comb{I})~\comb{K}) \betasteps~ \\
            \betasteps~ & \comb{K}
        \end{align}
    \end{proof}

    \subsection{Программирование в лямбда-исчислении}

    Локальный минимум: 3б.

    \subsubsection{(1б)}

    Задайте лямбда-терм $\textbf{or}~b_1~b_2$, который вычисляет дизъюнкцию двух логических величин.

    \begin{proof}
        \(\textbf{or} \coloneqq \lambda l~r\ldotp l~\textbf{true}~r\)

        Примеры:
        \begin{gather}
            \textbf{or}~\textbf{true}~\textbf{false} \betasteps\textbf{true}~\textbf{true}~\textbf{false}~\betastep(\lambda x\ldotp \textbf{true})~\textbf{false} \betastep\textbf{true} \\
            \textbf{or}~\textbf{false}~\textbf{true} \betasteps\textbf{false}~\textbf{true}~\textbf{false}~\betastep(\lambda x\ldotp x)~\textbf{true} \betastep\textbf{true}
        \end{gather}
    \end{proof}

    \subsubsection{(1б)}

    Задайте лямбда-терм $\term{mult}~a~b$, перемножающий два числа Чёрча, и лямбда-терм $\term{pow}~a~b$, который вычисляет $a^b$.

    \begin{proof}
        \(\textbf{mult} \coloneqq \lambda a~b\ldotp a~(\term{sum}~b)~\term{0}\)

        Примеры:
        \begin{gather}
            \term{mult}~\term{2}~\term{3} \betasteps \term{2}~(\term{sum}~\term{3})~\term{0} \betasteps \\
            \betasteps (\term{sum}~\term{3})~((\term{sum}~\term{3})~\term{0}) \betasteps \\
            \betasteps \term{sum}~\term{3}~\term{3} \betasteps \term{6}
        \end{gather}
        \begin{gather}
            \term{mult}~\term{0}~\term{2} \betasteps \term{0}~(\term{sum}~\term{2})~\term{0} \equiv \\
            \equiv (\lambda s~z\ldotp z)~(\term{sum}~\term{2})~\term{0} \betasteps \term{0}
        \end{gather}
        \begin{gather}
            \term{mult}~\term{1}~\term{0} \betasteps \term{1}~(\term{sum}~\term{0})~\term{0} \equiv \\
            \equiv (\lambda s~z\ldotp s~z)~(\term{sum}~\term{0})~\term{0} \betasteps \\
            \betasteps \term{sum}~\term{0}~\term{0} \betasteps \term{0}
        \end{gather}
    \end{proof}

    \subsubsection{(1б)}

    Задайте лямбда-терм $\term{isZero}$, который возвращает $\term{true}$, если подать $\term{0}$, и $\term{false}$ в противном случае.

    \begin{proof}
        \(\textbf{isZero} \coloneqq \lambda n\ldotp n~(\comb{K}~\term{false})~\term{true}\)

        Примеры:
        \begin{gather}
            \term{isZero}~\term{2} \betastep \term{2}~(\comb{K}~\term{false})~\term{true} \betasteps \\
            \betasteps (\comb{K}~\term{false})~((\comb{K}~\term{false})~\term{true}) \betastep \\
            \betastep (\lambda y\ldotp \term{false})~((\comb{K}~\term{false})~\term{true}) \betastep \term{false}
        \end{gather}
        \begin{gather}
            \term{isZero}~\term{0} \betastep \term{0}~(\comb{K}~\term{false})~\term{true} \betasteps \term{true}
        \end{gather}
    \end{proof}

    \subsubsection{(1б)}

    Скорее всего, ваш любимый язык программирования имеет поддержку функций высших порядков.
    Представьте, что в вашем языке нет структур, классов и т.д., и закодируйте на нём пары в стиле чистого лямбда-исчисления (\emph{не забывайте про каррирование}!).

    Как известно, котики полностью описываются двумя параметрами: именем и мурлычностью.
    Реализуйте библиотеку функций для работы с котиками, используя пары по Чёрчу как низкоуровневый примитив.
    Нужно написать функции: для создания котика, просмотра имени и мурлычности по котику.
    А так же реализовать функцию ``почесать'', которая по котику будет возвращать котика повышенной мурлычности.

    Если у вас статически типизированный язык программирования:
    \begin{itemize}
        \item Вам полагается дополнительный балл в случае успеха.
        \item Достаточно поддержать только хранение интов.
        \item Напишите вручную типы основных сущностей, для остального можно воспользоваться механизмом вывода типов вашего языка: \texttt{auto} в C++, \texttt{var} в Java\ldots
    \end{itemize}

    \begin{proof}
\begin{minted}[frame=lines,linenos=true,mathescape]{ocaml}
module Pair = struct
    (* Sorry for using struct, small hack 
       forall quantifier for return type only in fields *)
    type ('a, 'b) t = { f : 'r. ('a -> 'b -> 'r) -> 'r }

    let make_pair : 'a -> 'b -> ('a, 'b) t = fun a b -> { f = fun f -> f a b }
    let fst : ('a, 'b) t -> 'a = fun p -> p.f (fun a b -> a)
    let snd : ('a, 'b) t -> 'b = fun p -> p.f (fun a b -> b)
end

module Church = struct
    (* Same here *)
    type t = { f : 'a. ('a -> 'a) -> 'a -> 'a }

    let zero : t = { f = fun s z -> z }
    let succ : t -> t = fun n -> { f = (fun s z -> s @@ n.f s z) }

    let rec from_int = function
    | 0 -> zero
    | n -> succ @@ from_int @@ n - 1

    let to_int n = n.f Int.succ 0
end

module Cat : sig
    type t

    val make_cat : string -> int -> t
    val name : t -> string
    val meowness : t -> int
    val scratch : t -> t
end = struct
    open Pair
    open Church

    type t = (string, Church.t) Pair.t

    let make_cat_impl name meowness : t = 
        make_pair name @@ meowness

    let meowness_impl (cat : t) : Church.t = snd cat

    let make_cat name (meowness : int) : t = 
        make_cat_impl name @@ from_int meowness

    let name (cat : t) : string =
        fst cat

    let meowness (cat : t) : int =
        to_int @@ meowness_impl cat

    let scratch (cat) : t =
        let n = name cat in
        make_cat_impl n (succ @@ meowness_impl cat)
end

let main = 
    (* Check Church numerals *)
    assert Church.(2 = to_int @@ succ @@ succ @@ zero);
    assert Church.(10 = to_int @@ from_int 10);

    (* Check Pair *)
    assert Pair.(1 = fst @@ make_pair 1 2);
    assert Pair.(2 = snd @@ make_pair 1 2);

    (* Check cat *)
    let cat = Cat.make_cat "barsik" 2 in
    assert Cat.("barsik" = name cat);
    assert Cat.(2 = meowness cat);
    let cat = Cat.scratch cat in
    assert Cat.("barsik" = name cat);
    assert Cat.(3 = meowness cat);
\end{minted}
    \end{proof}

    \subsubsection{(1б)}

    Задайте лямбда-терм $\term{divides3}~n$, который возвращает $\term{true}$, если $n$ делится на 3 без остатка, и $\term{false}$ в противном случае.

    \begin{proof}
        Можно использользовать пары для создания троек, но получится громоздко, поэтому заведен отдельные тройки:
        \begin{align*}
            \term{mkTriple}  & \coloneqq \lambda x~y~z~f\ldotp f~x~y~z \\
            \term{rotTriple} & \coloneqq \lambda t\ldotp t~\lambda x~y~z~f\ldotp f~z~x~y \\
            \term{fstTriple} & \coloneqq \lambda t\ldotp t~\lambda x~y~z\ldotp x \\
            \term{divides3}   & \coloneqq \lambda n\ldotp \term{fstTriple}~(n~\term{rotTriple}~(\term{mkTriple}~\term{true}~\term{false}~\term{false}))
        \end{align}
        Примеры:
        \begin{gather}
            \term{divides3}~\term{2} \betasteps \term{fstTriple}~(\term{rotTriple}~(\term{rotTriple}~(\term{mkTriple}~\term{true}~\term{false}~\term{false}))) \betasteps \\
            \betasteps \term{fstTriple}~(\term{rotTriple}~(\term{rotTriple}~(\lambda f\ldotp f~\term{true}~\term{false}~\term{false}))) \betasteps \\
            \betasteps \term{fstTriple}~(\term{rotTriple}~((\lambda f\ldotp f~\term{true}~\term{false}~\term{false})~\lambda x~y~z~f\ldotp f~z~x~y)) \betasteps \\
            \betasteps \term{fstTriple}~(\term{rotTriple}~((\lambda x~y~z~f\ldotp f~z~x~y)~\term{true}~\term{false}~\term{false})) \betasteps \\
            \betasteps \term{fstTriple}~(\term{rotTriple}~(\lambda f\ldotp f~\term{false}~\term{true}~\term{false})) \betasteps \\
            \betasteps \term{fstTriple}~(\lambda f\ldotp f~\term{false}~\term{false}~\term{true}) \betasteps \\
            \betasteps (\lambda f\ldotp f~\term{false}~\term{false}~\term{true})~\lambda x~y~z\ldotp x \betasteps \\
            \betasteps (\lambda x~y~z\ldotp x)~\term{false}~\term{false}~\term{true} \betasteps \term{false}
        \end{gather}
        \begin{gather}
            \term{divides3}~\term{3} \betasteps \term{fstTriple}~(\term{rotTriple}~(\term{rotTriple}~(\term{rotTriple}~(\term{mkTriple}~\term{true}~\term{false}~\term{false})))) \betasteps \\
            \betasteps \term{fstTriple}~(\term{rotTriple}~(\term{rotTriple}~(\term{rotTriple}~(\lambda f\ldotp f~\term{true}~\term{false}~\term{false})))) \betasteps \\
            \betasteps \term{fstTriple}~(\term{rotTriple}~(\term{rotTriple}~(\lambda f\ldotp f~\term{false}~\term{true}~\term{false}))) \betasteps \\
            \betasteps \term{fstTriple}~(\term{rotTriple}~(\lambda f\ldotp f~\term{false}~\term{false}~\term{true})) \betasteps \\
            \betasteps \term{fstTriple}~(\lambda f\ldotp f~\term{true}~\term{false}~\term{false}) \betasteps \term{true}
        \end{gather}
    \end{proof}
         

    \subsubsection{(2б)}

    Задайте лямбда-терм $\term{pred}~n$, который находит предыдущее число Чёрча.

    Подсказка:
    \begin{lstlisting}[language=Python]
    def pred(n):
      prd = 0
      cur = 0
      for i in range(0, n):
        prd = cur
        cur = cur + 1
      return prd
    \end{lstlisting}

    \begin{proof}
        \begin{align*}
            \term{pred} \coloneqq \lambda n\ldotp \term{fst}~(n~(\lambda p\ldotp \term{pair}~(\term{snd}~p)~(\term{suc}~(\term{snd}~p)))~(\term{pair}~\term{0}~\term{0}))
        \end{align}
        Примеры:
        \begin{gather}
            \term{pred}~\term{2} \betasteps \term{fst}~((\lambda p\ldotp \term{pair}~(\term{snd}~p)~(\term{suc}~(\term{snd}~p)))~((\lambda p\ldotp \term{pair}~(\term{snd}~p)~(\term{suc}~(\term{snd}~p)))~(\term{pair}~\term{0}~\term{0}))) \betasteps \\
            \term{pred}~\term{2} \betasteps \term{fst}~((\lambda p\ldotp \term{pair}~(\term{snd}~p)~(\term{suc}~(\term{snd}~p)))~(\term{pair}~\term{0}~(\term{suc}~\term{0}))) \betasteps \\
            \term{pred}~\term{2} \betasteps \term{fst}~((\lambda p\ldotp \term{pair}~(\term{snd}~p)~(\term{suc}~(\term{snd}~p)))~(\term{pair}~\term{0}~\term{1})) \betasteps \\
            \term{pred}~\term{2} \betasteps \term{fst}~(\term{pair}~\term{1}~(\term{suc}~\term{1})) \betasteps \\
            \term{pred}~\term{2} \betasteps \term{fst}~(\term{pair}~\term{1}~\term{2}) \betasteps \term{1}
        \end{gather}
        \begin{gather}
            \term{pred}~\term{0} \betasteps \term{fst}~(\term{pair}~\term{0}~\term{0}) \betasteps \term{0}
        \end{gather}
    \end{proof}

\end{document}
